%============================================================ 4. 제안 방법
\section{제안 방법}
\label{sec:method}

% figure* 는 선언 다음 쪽에 놓인다. §4 제목 직후 선언 → §4 첫 인용 쪽(다음 쪽) 상단.
\begin{figure*}[!t]
\centering
\includegraphics[width=0.8\textwidth]{fig4.png}
\caption{판단 지연이 다른 2계층 비동기 구조}
\label{fig:architecture}
\end{figure*}

\subsection{2계층 비동기 구조}
\label{sec:arch}

제안 체계는 방어 판단을 전략 계층과 기동 계층으로 분리한다(Fig.~\ref{fig:architecture};
두 계층은 25초 주기로 결정하고 저수준 조타(PD)는 1초 주기로 돈다). 전략 계층은 \emph{어느 방어정이 어느 적 무리를 맡는가}를 정하고, 기동 계층은 \emph{배정된
무리를 막기 위해 해면 어디에 그물벽을 놓는가}를 정한다. 두 계층은 같은 주기로 결정을 내리지만
\textbf{응답 속도가 다르다}. 기동 계층의 순전파는 순간에 끝나는 반면, 전략 계층의
언어모델 호출은 응답에 지연이 발생한다(\S\ref{sec:res_latency}). 느린 전략 계층 호출만
비동기로 분리함으로써, 호출이 진행되는 동안 기동 계층은 직전 배정을 그대로 사용하여 결정을
계속하고 응답이 도착하면 다음 결정부터 새 배정이 반영된다.

\S\ref{sec:results}에서 지연은 설계 가정이 아니라 실측값으로 보고되며(\S\ref{sec:res_latency}), 비동기 분리는 그 실측이 요구한 결과다.

\subsection{적 무리 군집화}
\label{sec:cluster}

두 판단 계층은 개별 적선이 아니라 \emph{무리}를 단위로 배정한다. 따라서 매 결정마다 생존
적선을 모선 기준 방위로 묶어 최대 $C$ 개의 무리로 요약한다. 적선 $m$ 의 모선 기준 방위를
항법 규약(북쪽 $0^\circ$, 시계 방향)으로 정의한다. 절차는 Fig.~\ref{fig:clustering}의
네 단계 --- 모선 기준 방위각 정렬, 임계 $\delta_{\mathrm{gap}}$ 를 넘는 방위 빈틈 탐색,
빈틈에서 잘라 $1$--$C$개 무리 확정, 무리별 위협도와 요격점(무리 중심--모선 선분 위) 산출
--- 로 이루어진다.
\begin{equation}
\vartheta_m = \operatorname{atan2}\!\big(e_{m,x} - m_x,\; e_{m,y} - m_y\big) \bmod 360^\circ
\label{eq:bearing}
\end{equation}

\begin{figure}[!tbp]
\centering
\includegraphics[width=\columnwidth]{fig5.png}
\caption{방위각 정렬, 군집화, 군집 배정}
\label{fig:clustering}
\end{figure}

고정 방위 구간(bin)은 한 무리가 경계에 걸리면 둘로 쪼개지므로, \emph{방위 간격(gap)}으로
무리 수를 데이터에 맞추는 적응형 군집화를 쓴다. 생존 적선
$n = |\mathcal{A}|$ 척의 방위를 오름차순으로 정렬하여 $\vartheta_{(1)} \le \dots \le
\vartheta_{(n)}$ 이라 하고, 이웃한 방위 사이의 간격을 원형으로 정의한다.
\begin{equationfn}
g_{(q)} = \vartheta_{(q+1)} - \vartheta_{(q)}\ (q < n), \quad
g_{(n)} = \vartheta_{(1)} + 360^\circ - \vartheta_{(n)}
\label{eq:gap}
\end{equationfn}
임계 $\delta_{\mathrm{gap}}$ 보다 큰 간격의 수가 무리 수를 정하되, $1$ 이상 $C$ 이하(그리고
적선 수 이하)로 자른다.
\begin{equation}
\hat{C} = \mathrm{clip}\!\Big(\big|\{\, q : g_{(q)} > \delta_{\mathrm{gap}} \,\}\big|,\; 1,\; \min(C, n)\Big)
\label{eq:ncluster}
\end{equation}
가장 큰 간격 $\hat{C}$ 개를 무리 경계로 삼고, 그 가운데 최대 간격을 원형 순서의 이음새로
열어 이음새 다음 적선부터 방위 순으로 걸으며 경계를 지날 때마다 무리 번호를 하나씩
올린다. 이렇게 얻은 라벨 $j_m \in \{1, \dots, \hat{C}\}$ 로 무리 $j$ 의 구성원 집합
$\mathcal{A}_j = \{m : j_m = j\}$ 가 정해진다. 무리 $j$ 는 네 가지 양으로 요약한다.
\begin{equation}
\begin{aligned}
n_j &= |\mathcal{A}_j|, \qquad
\mathbf{c}_j = \frac{1}{n_j} \sum_{m \in \mathcal{A}_j} \mathbf{e}_m,\\
\Delta\vartheta_j &= 90^\circ \sqrt{1 - \bar{R}_j}, \qquad
\bar{R}_j = \Big\lVert \frac{1}{n_j} \sum_{m \in \mathcal{A}_j} \hat{\mathbf{b}}_m \Big\rVert,\\[2pt]
v^{\mathrm{app}}_j &= \frac{v_e}{n_j} \sum_{m \in \mathcal{A}_j}
  \big\langle (\sin\psi_m, \cos\psi_m),\, \hat{\mathbf{d}}_m \big\rangle,\\
\hat{\mathbf{d}}_m &= (\mathbf{m} - \mathbf{e}_m) / \lVert \mathbf{m} - \mathbf{e}_m \rVert
\end{aligned}
\label{eq:clustersum}
\end{equation}
$n_j$ 는 척수, $\mathbf{c}_j$ 는 무리 중심, $\Delta\vartheta_j$ 는 방위 단위벡터
$\hat{\mathbf{b}}_m = (\cos\vartheta_m, \sin\vartheta_m)$ 의 원형 결집도 $\bar{R}_j$ 에서 얻은
각도 퍼짐(한 점에 모이면 $0^\circ$), $v^{\mathrm{app}}_j$ 는 모선
방향 평균 접근 속도이다. 본 연구는 $C = 4$, $\delta_{\mathrm{gap}} \approx 12^\circ$ 를 쓴다.
이 규칙 아래에서 집중 포메이션은 무리 하나, 양동 포메이션은 세 방위의 무리 셋으로
요약된다. 파상 포메이션은 단(rank)마다 방위를 $18^\circ$ 씩 벌려 생성되므로
$\delta_{\mathrm{gap}}$ 을 넘어 대개 단마다 별개의 무리가 되며, 단 사이의 거리차는 무리
중심의 반경 차이와 \S\ref{sec:obs}의 격자 관측에 함께 남는다. 척수와 중심은
\S\ref{sec:heuristic}의 위협도와 요격점을, 네 요약량 전부는 \S\ref{sec:commander}의
지휘관 입력을 구성한다.

\subsection{기하 휴리스틱 기준 방법}
\label{sec:heuristic}

\S\ref{sec:cluster}의 무리 요약 위에서 같은 문제를 기하 규칙만으로 푸는 기준 방법으로,
\S\ref{sec:results}의 비교 기준이자 \S\ref{sec:bc}의 행동 복제 목표다.

\textbf{배정.} 위협도 $\tau_j$ 는 척수와 모선 근접도의 곱, 요격점 $\mathbf{I}_j$ 는 무리
중심--모선 선분 위에서 모선 쪽으로 치우친 점이다($R_w$: 해역 반폭).
\begin{equation}
\begin{aligned}
\tau_j &= n_j \cdot \mathrm{clip}\!\left(1 - \frac{\lVert \mathbf{c}_j - \mathbf{m}\rVert}{R_w},
         \, 0.05,\, 1\right),\\
\mathbf{I}_j &= \mathbf{c}_j + \lambda_{\mathrm{I}}\,(\mathbf{m} - \mathbf{c}_j), \qquad \lambda_{\mathrm{I}} = 0.55
\end{aligned}
\label{eq:threat}
\end{equation}
위협 상위 $P$ 개 무리에 대해, 요격점 거리 $d_{pj} = \lVert \mathbf{a}_p - \mathbf{I}_j \rVert$
에서 직전 배정 유지 쌍에 $\beta$ 를 감한 비용으로 탐욕적 1:1 배정을 한다(전개 중 담당이
바뀌면 이동과 그물이 회수되지 않기 때문).
\begin{equation}
\tilde{d}_{pj} = d_{pj} - \beta \cdot \mathbb{1}[\,j = j_p^{\text{prev}}\,]
\label{eq:sticky}
\end{equation}
활성 무리가 $P$ 보다 적으면 남는 방어정은 부하가 가장 큰 무리에 안쪽 층($\lambda_{\mathrm{L}}
= 0.18$, 방위 $\pm 40^\circ$)으로 덧붙인다.

\textbf{경유점.} 요격점을 가운데 둔, 접근 방향에 수직인 길이 $L_{\mathrm{net}} = 450\unit{m}$
의 벽이므로 $K = 2$ 개다($\hat{\mathbf{r}}$: 모선$\to$요격점 단위 벡터).
\begin{equationsmall}
\mathbf{t}_k = \mathbf{I} + \left(k - \tfrac{K-1}{2}\right) L_{\mathrm{net}}\,
               \hat{\mathbf{r}}^{\perp}, \qquad k = 0,\dots,K-1
\label{eq:heur_target}
\end{equationsmall}
각 $\mathbf{t}_k$ 는 학습 정책과 같은 격자·유효 마스크·간격 제약(\S\ref{sec:validmask}) 안의
최근접 픽셀로 스냅하므로 두 방법의 행동공간은 같다. 휴리스틱은 기하로 정해진 한 점만
지목하고 다른 방어정의 벽, 파상 공격의 뒤 단, 지형의 길목을 보지 못한다 --- 이 해면 전체의
비교를 학습된 점수맵(\S\ref{sec:policy})이 대체하므로, \S\ref{sec:res_maneuver}의 차이는
표현력이 아니라 \emph{결정 규칙}의 차이다.

\subsection{점수맵 기반 기동 정책}
\label{sec:policy}

\subsubsection{다채널 격자 관측}
\label{sec:obs}

기동 계층의 입력은 좌표 목록이 아니라, 모선을 중앙에 두고 관측 반폭 $E$ 의 전장을
$n_{\mathrm{g}} \times n_{\mathrm{g}}$ 격자로 나눈 다채널 지도이다(채널별 실제 화면은
Fig.~\ref{fig:observation}). 총 15개 채널은 모든 방어정이 공유하는 전역 9채널, 방어정마다
다른 3채널, CoordConv~\citep{liu2018coordconv} 좌표 3채널로 구성되며
(Fig.~\ref{fig:tensorstack}), 각 채널의 정의와 값 범위를 Table~\ref{tab:channels}에
정리하였다. 그림의 좌측 시선원뿔과 함정 형상은 삽화의 표현이다 --- 본 연구의 래스터는 방어정 시야가
아니라 모선을 중앙에 둔 전역 격자이며, 방어정마다 달라지는 것은 3채널(10--12)과
$\mathbf{o}_p$ 뿐이다.

\begin{figure*}[tp]
\centering
\includegraphics[width=0.7\textwidth]{fig6.png}
\caption{래스터화된 15채널 관측의 예(파상 포메이션, 통영 매물도 북측)}
\label{fig:observation}
\end{figure*}
\begin{figure}[!htbp]
\centering
\includegraphics[width=\columnwidth]{fig7.png}
\caption{15채널 관측의 구성 --- 전역·방어정별·좌표 채널의 격자 중첩과 격자 밖 자체 상태 벡터}
\label{fig:tensorstack}
\end{figure}

\begin{table*}[!tp]
\centering
\caption{15채널 격자 관측의 구성. 전역 채널은 월드당 한 번 계산하여 모든 방어정이 공유하고,
         방어정별 채널만 방어정마다 다르다. 정의 열의 식 번호는 본문의 수식이다.}
\label{tab:channels}
\footnotesize
\renewcommand{\arraystretch}{1.15}
\begin{tabular}{@{}clllc@{}}
\toprule
\# & 묶음 & 채널 & 정의 & 범위 \\
\midrule
1  & \multirow{9}{*}{전역}   & 적 밀도        & 픽셀 내 생존 적선 수$/M$, 식~\eqref{eq:ch_enemy}      & $[0,1]$ \\
2  &                         & 적 진행 $x$    & 픽셀 내 적선 헤딩 단위벡터의 평균, 식~\eqref{eq:ch_enemy} & $[-1,1]$ \\
3  &                         & 적 진행 $y$    & 위와 같음                                              & $[-1,1]$ \\
4  &                         & 적 위협도      & 모선 도달 임박도의 픽셀 내 최대, 식~\eqref{eq:ch_threat} & $[0,1]$ \\
5  &                         & 아군 위치      & 픽셀 내 생존 방어정 수$/P$                             & $[0,1]$ \\
6  &                         & 설치된 그물    & 물리 격자의 그물 점유를 최대값 풀링                     & $\{0,1\}$ \\
7  &                         & 원점           & 모선 반경 안의 픽셀                                    & $\{0,1\}$ \\
8  &                         & 요격 구역      & $r_{\min} \le \lVert \mathbf{x}_i - \mathbf{m} \rVert \le r_{\max}$ 인 픽셀 & $\{0,1\}$ \\
9  &                         & 육지           & 고해상도 육지 마스크의 면적비  & $[0,1]$ \\
\midrule
10 & \multirow{3}{*}{방어정별} & 자기 위치     & 방어정 $p$ 가 놓인 픽셀                                & $\{0,1\}$ \\
11 &                         & 배정 요격점    & $\mathbf{I}_p$ 가 놓인 픽셀 (미배정이면 전부 0)        & $\{0,1\}$ \\
12 &                         & 유효 영역      & 유효 마스크 $\mathcal{V}_p$, 식~\eqref{eq:validset}    & $\{0,1\}$ \\
\midrule
13 & \multirow{3}{*}{좌표}   & $x$            & $(x_i - m_x)/E$, 식~\eqref{eq:ch_coord}                & $[-1,1]$ \\
14 &                         & $y$            & $(y_i - m_y)/E$                                        & $[-1,1]$ \\
15 &                         & $r$            & $\mathrm{clip}(\lVert \mathbf{x}_i - \mathbf{m} \rVert / E, 0, 1)$ & $[0,1]$ \\
\bottomrule
\end{tabular}
\end{table*}

\textbf{래스터화와 정규화.} 픽셀 크기는 $\ell = 2E/n_{\mathrm{g}}$ 이고, 월드 좌표
$\mathbf{x}$ 는 모선 위치 $\mathbf{m}$ 을 중심으로 다음 픽셀 인덱스로 사상된다.
\begin{equation}
\iota(\mathbf{x}) = \mathrm{clip}\!\left(\left\lfloor \frac{\mathbf{x} - \mathbf{m} + E\mathbf{1}}{\ell} \right\rfloor,\, 0,\, n_{\mathrm{g}}-1\right)
\label{eq:world2pix}
\end{equation}
적 관련 채널은 개체를 픽셀에 ``칠하는'' 방식으로 만들며, 값 범위가 $[0,1]$ 또는 $[-1,1]$
에 놓이도록 정규화한다. 생존 적선 집합을 $\mathcal{A}$, 적선 $m$ 의 위치·헤딩을
$\mathbf{e}_m$, $\psi_m$ 이라 하면 픽셀 $i$ 의 적 밀도·진행 방향·위협도 채널은
\begin{equationsmall}
c^{\mathrm{den}}_i = \frac{N_i}{M}, \quad
c^{\mathrm{vel}}_i = \frac{1}{\max(N_i, 1)} \sum_{m \in \mathcal{A}_i} (\sin\psi_m, \cos\psi_m)
\label{eq:ch_enemy}
\end{equationsmall}
\begin{equation}
c^{\mathrm{thr}}_i = \max_{m \in \mathcal{A}_i}
   \left(1 - \mathrm{clip}\!\left(\frac{\lVert \mathbf{e}_m - \mathbf{m} \rVert}{v_e T_{\max}},\, 0,\, 1\right)\right)
\label{eq:ch_threat}
\end{equation}
이다. 여기서 $\mathcal{A}_i = \{m \in \mathcal{A} : \iota(\mathbf{e}_m) = i\}$ 는 픽셀 $i$ 에
놓인 생존 적선의 집합이고 $N_i = |\mathcal{A}_i|$ 이다. 밀도는 적 총수 $M$ 으로, 진행 방향은 픽셀 안의 적 수로 나누어 평균 단위벡터가 되게
하고, 위협도는 모선까지 거리를 적이 에피소드 상한 안에 달릴 수 있는 최대 거리
$v_e T_{\max}$ 로 정규화하여 가까울수록 1에 가깝게 한다.
아군 위치 채널은 같은 방식으로 $P$ 로 나누고, 설치된 그물 채널은 고해상도 물리 격자의
최대값 풀링, 육지 채널은 면적비 다운샘플이다. 좌표 채널은
픽셀 중심 $\mathbf{x}_i = (x_i, y_i)$ 에 대해
\begin{equationsmall}
(c^{x}_i,\, c^{y}_i) = \frac{\mathbf{x}_i - \mathbf{m}}{E}, \qquad
c^{r}_i = \mathrm{clip}\!\left(\frac{\lVert \mathbf{x}_i - \mathbf{m} \rVert}{E},\, 0,\, 1\right)
\label{eq:ch_coord}
\end{equationsmall}
로 두어 $[-1,1]$ 과 $[0,1]$ 에 놓는다. 격자에 담지 않는 방어정 $p$ 의 자체 상태 벡터는
같은 반폭 $E$ 로 정규화한 위치와 요격점, 헤딩 단위벡터, 그물 잔량 비율, 전개 여부이다.
\begin{equation}
\begin{aligned}
\mathbf{o}_p = \Big[&\tfrac{1}{E}(\mathbf{a}_p - \mathbf{m});\; (\sin\psi_p, \cos\psi_p);\;
  \tfrac{n_p}{n_{\max}};\; \delta_p;\\
  &\mathbb{1}[\text{배정}]\,\tfrac{1}{E}(\mathbf{I}_p - \mathbf{m})\Big] \in \mathbb{R}^{8}
\end{aligned}
\label{eq:own}
\end{equation}
여기서 $n_p$ 는 잔여 그물 수, $n_{\max}$ 는 척당 보유량, $\delta_p \in \{0,1\}$ 은 전개 중
여부이며, 배정이 없으면 요격점 성분은 0이다.

\textbf{스케일 불변성.} 길이는 $E$ 나 $\ell$ 로, 개수는 $M$·$P$·$n_{\max}$ 로, 시간은
$v_e T_{\max}$ 로 나뉘므로 모든 입력은 $[-1,1]$ 의 무차원 양이다. 해역 크기나 격자 해상도를
바꾸어도 채널의 의미와 값 범위가 변하지 않으며, 통계 기반 정규화 층이 필요 없다($E$,
$n_{\mathrm{g}}$ 등의 값은 Table~\ref{tab:scenario}).

목록 대신 격자 지도를 쓰면 입력 형태가 개체 수와 순서에 불변이고, 위의 정규화와 결합하면
해역 스케일에도 불변이다. 합성곱은 평행이동 등변이라 절대 위치가 흐려지므로 좌표 채널을
함께 넣어 위치를 보존한다.

육지 채널은 OpenStreetMap 해안선·섬 폴리곤을 래스터화한 실해역 마스크이며(학습에 쓴 6곳은
Fig.~\ref{fig:sites}), 충돌 판정은 고해상도 물리 격자에서 이진으로, 관측 채널은 이를
면적비로 다운샘플한 연속값으로 쓴다.

\begin{figure}[!tbp]
\centering
\includegraphics[width=\columnwidth]{fig8.png}
\caption{학습에 쓴 실해역 6곳의 육지 마스크(별표 모선, 파선 원 요격 구역)}
\label{fig:sites}
\end{figure}

%-------------------------------------------------- 그림 4

\subsubsection{FiLM 조건화 U-Net과 점수맵}

\begin{figure*}[!tbp]
\centering
\includegraphics[width=0.75\textwidth]{fig9.png}
\caption{점수맵 정책의 구조 --- 방어정 상태 질의와 U-Net 특징맵의 내적으로 얻는 점수맵과 부화소 오프셋}
\label{fig:policy}
\end{figure*}

관측은 U-Net 구조~\citep{ronneberger2015unet}를 지나 입력과 같은 해상도의 특징 지도로
변환된다(Fig.~\ref{fig:policy}). 인코더 2단($50\to25\to13$)으로 수용 영역을 넓히고, 디코더
2단으로 원해상도를 복원하며, skip 연결이 인코더의 위치 정보를 디코더로 전달한다. 정밀 좌표와
전역 상황 판단을 동시에 요구하는 본 문제에 U자형 구조가 맞물리는 지점이 여기다.

병목에서는 전역 평균 풀링으로 압축한 전황 벡터로부터 FiLM~\citep{perez2018film} 의
채널별 배율·이동 계수 $(\boldsymbol{\gamma}_{\mathrm{F}}, \boldsymbol{\beta}_{\mathrm{F}})$ 를
만들어 병목 특징을 변조한 뒤 디코더에 넘긴다(아래첨자 F 는 FiLM 을 뜻하며, 식~\eqref{eq:sticky}의 $\beta$ 나
할인 계수 $\gamma$ 와는 무관하다). 병목 특징 지도를 $\mathbf{H}_b \in \mathbb{R}^{2w \times h \times h}$
($w = 32$, $h = 13$)라 하면
\begin{equation}
\begin{aligned}
(\boldsymbol{\gamma}_{\mathrm{F}}, \boldsymbol{\beta}_{\mathrm{F}})
  &= \mathbf{W}_{\mathrm{F}}\, \mathrm{GAP}(\mathbf{H}_b) + \mathbf{b}_{\mathrm{F}},\\
\tilde{\mathbf{H}}_b &= (1 + \boldsymbol{\gamma}_{\mathrm{F}}) \odot \mathbf{H}_b + \boldsymbol{\beta}_{\mathrm{F}}
\end{aligned}
\label{eq:film}
\end{equation}
이며, $\odot$ 은 채널별 곱이다. 배율을 $1 + \boldsymbol{\gamma}_{\mathrm{F}}$ 로 두어 초기에는
항등 변조에서 출발한다. 국소 합성곱만으로는 ``지금 전황이 급박한가''와 같은 전역 정보가
모든 위치에 고르게 전달되지 않기 때문에 이 경로를 둔다.

각 방어정 $p$ 는 자신의 상태 벡터 $\mathbf{o}_p \in \mathbb{R}^8$(식~\eqref{eq:own})로부터
질의 벡터 $\mathbf{q}_p = \mathrm{MLP}_{8 \to 128 \to 32}(\mathbf{o}_p) \in \mathbb{R}^{32}$ 를
만들고, 픽셀 $i$ 의 특징 $\mathbf{f}_i \in \mathbb{R}^{32}$ 와의 내적으로 점수를 얻는다.
\begin{equation}
s_{p,i} = \frac{\langle \mathbf{W}_q \mathbf{q}_p + \mathbf{b}_q,\, \mathbf{f}_i \rangle}{\sqrt{d}}, \qquad d = 32
\label{eq:score}
\end{equation}
여기서 $\mathbf{W}_q \in \mathbb{R}^{d \times d}$, $\mathbf{b}_q \in \mathbb{R}^{d}$ 는 질의를 특징 공간에 맞추는 아핀 사영이다.
이는 어텐션~\citep{vaswani2017attention}의 점수 계산과 같은 형식이며, 포인터
방식~\citep{vinyals2015pointer}으로 격자 위의 원소를 지목하는 구조에 해당한다. 특징 지도는
방어정마다 다르므로(15채널 중 3채널이 방어정별이다), 관측 텐서는 방어정을 배치 축으로
쌓아 \textbf{결정당 한 번의 배치 호출}로 U-Net을 통과한다. 합성곱 연산량은 방어정 수에
비례하지만, \textbf{가중치는 모든 방어정이 공유하므로 모델 크기는 방어정 수에 불변}이다.
방어정별로 달라지는 것은 질의 $\mathbf{q}_p$ 와 방어정별 3채널뿐이다.

합성곱 블록은 $3\times3$ 합성곱 -- 그룹 정규화(8그룹)~\citep{wu2018groupnorm} -- SiLU 다.
그룹 정규화는 표본 하나 안에서 정규화하므로 running statistics 가 없고, 배치 크기 1의
온보드 추론이 학습과 같은 연산을 수행한다. 같은 배치 관점에서 정책은 조타 대신 경유점만 내어
선체 동역학~\citep{fossen2011marine}과 분리되고, 관측·행동이 격자 상대 좌표라 실해역
이식 시 축척만 재설정하면 된다 --- 세 성질은 시뮬레이션 안에서만 확인되었다(\S\ref{sec:limitations}).

\subsubsection{유효 마스크와 픽셀 지목}
\label{sec:validmask}

점수맵 전체가 후보가 되는 것은 아니다. 방어정 $p$ 의 유효 픽셀 집합 $\mathcal{V}_p$ 는 네
조건 --- 환상 요격 구역 안, 육지 밖, 배정 요격점의 반경 게이트 안, 배정 무리의 방위 회랑
안 --- 의 교집합이다. 픽셀 $i$ 의 중심을 $\mathbf{x}_i$, 육지 픽셀 집합을
$\mathcal{L}_{\mathrm{land}}$, 배정 요격점을 $\mathbf{I}_p$, 모선에서 본 요격점 방위와 픽셀
방위를 $\vartheta_p$, $\vartheta_i$ 라 하면
\begin{equation}
\begin{aligned}
\mathcal{V}_p = \big\{\, i :\; & r_{\min} \le \lVert \mathbf{x}_i - \mathbf{m} \rVert \le r_{\max},\;
  i \notin \mathcal{L}_{\mathrm{land}},\\
  & \lVert \mathbf{x}_i - \mathbf{I}_p \rVert \le r_{\mathrm{g}},\;
  |\vartheta_i - \vartheta_p|_{\circ} \le \vartheta_{\mathrm{g}} \,\big\}
\end{aligned}
\label{eq:validset}
\end{equation}
여기서 $|\cdot|_{\circ}$ 는 $[0^\circ, 180^\circ]$ 로 접은 각도 차이이며 $r_{\min} = 400$,
$r_{\max} = 4500$, $r_{\mathrm{g}} = 3000\unit{m}$, $\vartheta_{\mathrm{g}} = 60^\circ$ 이다.
유효 픽셀이 하한 150개에 못 미치면 반경 게이트부터 푼다 --- 방위 회랑은 ``무엇을 막는
그물인가''를 정하므로 마지막까지 유지한다. 마스크는 로짓 단계에서 적용된다.
\begin{equationsmall}
\tilde{s}_{p,i} = \begin{cases} s_{p,i}, & i \in \mathcal{V}_p \\ -\infty, & i \notin \mathcal{V}_p, \end{cases}
\quad
\pi_\theta(i \mid \mathbf{s}, p) = \frac{e^{\tilde{s}_{p,i}}}{\sum_{j \in \mathcal{V}_p} e^{\tilde{s}_{p,j}}}
\label{eq:masked_softmax}
\end{equationsmall}
유효 집합 밖 픽셀의 확률은 정확히 0 이다. 즉 안전·임무 제약이 손실 항으로 ``되도록
지키게'' 유도되는 것이 아니라 \emph{행동 분포에 직접} 부과되어, 제약 위반이 애초에 표현
불가능하다. 이 제약은 강하다 --- 학습된 체크포인트로 잰 유효 비율은 격자의 평균 10.5\,\%
에 그친다(Fig.~\ref{fig:validmask}(c)). 격자 대부분이 애초에 갈 수 없는 자리이며, 연속 좌표
회귀에는 이 구조를 표현할 수단이 없다.

\begin{figure*}[tbp]
\centering
\includegraphics[width=\textwidth]{fig10.png}
\caption{유효 마스크와 픽셀 지목(학습된 체크포인트의 실제 출력): (a) 유효 영역, (b) 점수맵과 선택된 픽셀 쌍,
         (c) 제약별 후보 수 감소 --- 대표 사례에서 $2500 \to 1000 \to 979 \to 387 \to 299$ 칸;
         6{,}652회 결정 평균 유효 비율 10.5\,\%}
\label{fig:validmask}
\end{figure*}

이산 격자는 마스킹 외에 두 이득을 준다. 범주형 표본은 유효 영역을 벗어나지 않고 엔트로피가
유한하므로 \S\ref{sec:training}의 그룹 후보들이 실제로 서로 다른 칸을 밟고(가우시안 정책은
표준편차가 줄면 후보가 한 점으로 뭉친다), 관측 격자와 행동 격자가 같아 합성곱의 평행이동
등변성이 ``이 칸이 좋다''는 판단으로 곧바로 이어진다. 한 칸 252\,m 의 해상도는 아래의 칸 안
연속 오프셋으로 보완한다.

각 방어정은 유효 영역 안에서 점수 상위 픽셀 $K=2$ 개를 \emph{순차적으로} 고른다. $k$ 번째
선택의 질의와 후보 집합을 $\mathbf{q}^{(k)}_p$, $\mathcal{V}^{(k)}_p$ 라 하면
($\mathbf{q}^{(0)}_p = \mathbf{q}_p$, $\mathcal{V}^{(0)}_p = \mathcal{V}_p$)
\begin{equationfn}
i_k \sim \pi^{(k)}_p \equiv \pi_\theta(\cdot \mid \mathbf{q}^{(k)}_p, \mathcal{V}^{(k)}_p), \qquad
\mathbf{q}^{(k+1)}_p = \mathbf{q}^{(k)}_p + \mathbf{f}_{i_k}
\label{eq:autoreg}
\end{equationfn}
\begin{equationsmall}
\mathcal{V}^{(k+1)}_p = \big\{ i \in \mathcal{V}^{(k)}_p :
  r_{\mathrm{dup}} \le \lVert \mathbf{x}_i - \mathbf{x}_{i_k} \rVert \le 1.5\, L_{\mathrm{net}} \big\}
\label{eq:exclude}
\end{equationsmall}
이다. 먼저 고른 픽셀의 특징을 질의에 더해 두 번째를 고르므로 두 점이 서로를 알고 뽑히며,
식~\eqref{eq:exclude}는 두 점 사이 간격의 상·하한이다 --- $r_{\mathrm{dup}} = 200\unit{m}$
안쪽은 벽이 서지 않고, 상한 $1.5\,L_{\mathrm{net}} = 675\,$m 는 칠 수 없는 길이의 벽을
막는다. 제약으로 후보가 모두 사라지면 제약을 완화하여 결정을 반환한다.

선택된 칸 안의 위치는 연속 오프셋으로 정한다. 질의와 선택 픽셀의 특징으로부터 평균을 내고,
학습되는 로그표준편차 $\boldsymbol{\sigma}_o$ 를 가진 가우시안에서 표집한다.
\begin{equationsmall}
\begin{aligned}
\boldsymbol{\mu}_k &= \mathrm{MLP}_{2d \to d \to 2}\big([\mathbf{q}^{(0)}_p;\, \mathbf{f}_{i_k}]\big), \quad
\mathbf{u}_k \sim \mathcal{N}(\boldsymbol{\mu}_k, \boldsymbol{\sigma}_o^2),\\
\mathbf{t}_k &= \mathbf{x}_{i_k} + \tfrac{\ell}{2}\,\mathrm{clip}(\mathbf{u}_k, -1, 1)
\end{aligned}
\label{eq:offset}
\end{equationsmall}
즉 \emph{어느 칸인가}는 범주형으로, \emph{칸 안 어디인가}는 가우시안으로 결정된다. 한
방어정의 행동 $a_p = \{(i_k, \mathbf{u}_k)\}_{k=0}^{K-1}$ 의 로그확률은 두 분포의 합이며,
엔트로피 $\mathcal{H}(a_p)$ 도 같은 꼴로 정의된다.
\begin{equationsmall}
\log \pi_\theta(a_p) = \sum_{k=0}^{K-1} \Big[ \log \pi^{(k)}_p(i_k)
  + \log \mathcal{N}(\mathbf{u}_k; \boldsymbol{\mu}_k, \boldsymbol{\sigma}_o^2) \Big]
\label{eq:logp}
\end{equationsmall}
특징 조회에서 두 번째 점까지의 흐름을 Fig.~\ref{fig:picksteps}에 요약하였다.

\begin{figure*}[!tbp]
\centering
\includegraphics[width=0.9\textwidth]{fig11.png}
\caption{점수맵 위 픽셀 지목의 네 단계 --- 픽셀 특징, 유효 마스크, 첫 점 선택, 간격 제약 아래 둘째 점 선택}
\label{fig:picksteps}
\end{figure*}

선택된 두 좌표를 잇는 선분이 그물벽 경로가 된다. 방어정은 가까운 점까지 무전개로 이동한 뒤
다음 구간에서 그물을 전개한다. 운용 시에는 최대 점수를 취하고(greedy), 학습 시에는 확률
분포에서 표본을 뽑는다.

\textbf{정지 규칙.} 그물이 없거나 배정되지 않은 방어정은 정지시킨다. 이런 방어정은 학습
중 보상 신호가 없어 점수맵이 거의 평평하므로(최고 확률 0.02, 배정된 방어정은 $0.5\sim0.7$),
그 출력이 항해 명령으로 새어 나가지 않게 한다.

\subsection{가치망 없는 다중 에이전트 그룹 상대 최적화}
\label{sec:training}

\subsubsection{그룹 상대 이득}

\begin{figure*}[tp]
\centering
\includegraphics[width=0.6\textwidth]{fig12.png}
\caption{가치망 없는 그룹 상대 최적화의 한 갱신 --- 후보 표집(위), 롤아웃 창 평가(가운데), 상대 비교와
         반사실적 크레딧(아래)}
\label{fig:grpocredit}
\end{figure*}

학습은 가치함수를 따로 두지 않는 그룹 상대 정책 최적화~\citep{shao2024deepseekmath}를 다중
방어정 상황에 맞게 변형하여 수행한다. 결정 시점의 상태를 $\mathbf{s}$, 방어정 $P$척의 결합
행동을 $\mathbf{a} = (a_1, \dots, a_P)$ 라 하자. 매 결정 시점에서 정책으로부터 $G$개의 결합
행동 후보를 표집하되, 후보 가운데 하나($k=0$)는 휴리스틱(\S\ref{sec:heuristic})이 내는
행동으로 고정하여 기준선으로 삼는다.
\begin{equationfn}
\mathbf{a}^{(0)} = \mathbf{a}_{\mathrm{heur}}(\mathbf{s}), \qquad
\mathbf{a}^{(k)} \sim \pi_\theta(\cdot \mid \mathbf{s}), \quad k = 1, \dots, G-1
\label{eq:candidates}
\end{equationfn}
각 후보를 시뮬레이터에서 공통 구간 $H = 180\,\mathrm{step}$ 만큼 실제로 진행시켜 창 보상
$R_k = R(\mathbf{s}, \mathbf{a}^{(k)}; H)$ 를 얻은 뒤(\S\ref{sec:reward}), 후보들 사이의
상대 비교로 이득을 산정한다.
\begin{equation}
A_k = \frac{R_k - R_0}{\mathrm{std}(\{R_j\}_{j=0}^{G-1}) + \epsilon}, \qquad \epsilon = 10^{-6}
\label{eq:adv}
\end{equation}
$R_0 = R_{\mathrm{heur}}$ 이므로 휴리스틱과 같은 성과를 내는 후보의 이득은 0 이고, 휴리스틱을
넘어서는 후보만 양의 이득을 받는다. 그룹 표준편차로 나누므로 이득은 크기가 없는 양이며,
그룹 안의 순위 정보만 남는다. 어떤 선택이 더 나았는지를 실제 결과로 직접 비교하므로 별도의
가치망이 필요하지 않다. 이는 가치 추정 편향과 비평가 학습의 불안정성을 원천적으로 제거한다.

한 갱신은 $N$개의 병렬 전장(월드)에서 동시에 수행된다. 월드 $n$ 의 후보 보상이 모두 같으면
이득이 정의되지 않으므로, 유효 마스크 $w_n = \mathbb{1}[\mathrm{std}(\{R_{n,k}\}_k) > \epsilon_{\mathrm{var}}]$
($\epsilon_{\mathrm{var}} = 10^{-6}$)로 그런 월드를 갱신에서 제외한다. 정책경사 손실은 유효 월드의 모든 후보·모든 방어정에 대해 이득으로 가중한 로그확률
(식~\eqref{eq:logp})의 음수이다.
\begin{equationfn}
\mathcal{L}_{\mathrm{PG}}(\theta) = -\frac{\sum_{n} w_n \sum_{k=0}^{G-1} A_{n,k} \sum_{p=1}^{P} \log \pi_\theta\big(a^{(k)}_{n,p} \mid \mathbf{s}_n\big)}{G \sum_n w_n}
\label{eq:lpg}
\end{equationfn}
여기에 엔트로피 정칙화를 더한 총손실을 최소화한다.
\begin{equationsmall}
\mathcal{L}(\theta) = \mathcal{L}_{\mathrm{PG}}(\theta)
  - \lambda_{\mathrm{H}}(t)\, \frac{1}{NP} \sum_{k=0}^{G-1} \sum_{n=1}^{N} \sum_{p=1}^{P} \mathcal{H}\big(a^{(k)}_{n,p}\big)
\label{eq:loss}
\end{equationsmall}
엔트로피는 $G$개 후보 모두에 대해 합하므로 유효 계수는 $G \lambda_{\mathrm{H}}(t)$ 이며,
$\lambda_{\mathrm{H}}(t)$ 는 \S\ref{sec:schedule}의 감쇠 계수이다. 기울기 노름은
$g_{\max} = 1.0$ 에서 잘라 Adam 으로 갱신한다. 이 절차와 다음 소절의 크레딧 배분을
Fig.~\ref{fig:grpocredit}에 정리하였다(그림은 $G=4$, $H=4$ 로 줄여 그렸으며 실제 학습은
$G=8$, $H=180$ step 이다).

\subsubsection{반사실적 크레딧 배분: 비교 변형}

식~\eqref{eq:lpg}는 팀 보상의 이득을 모든 방어정에 같은 값으로 곱하므로 개별 기여가 가려질
수 있다. 이를 분리하는 변형으로, 기준 표본 $\mathbf{a}^{\mathrm{base}} \sim
\pi_\theta(\cdot \mid \mathbf{s})$ 에서 방어정 $p$ 의 행동만 바꾼 혼합 후보 $G$개(후보 0 은
휴리스틱)를 같은 구간 $H$ 만큼 재진행하여 창 보상 $R^{(p)}_k$ 를 얻고, 방어정 $p$ 만의 이득
\begin{equation}
A^{(p)}_k = \frac{R^{(p)}_k - R^{(p)}_0}{\mathrm{std}(\{R^{(p)}_j\}_{j=0}^{G-1}) + \epsilon}
\label{eq:cfadv}
\end{equation}
을 방어정 $p$ 의 로그확률에만 곱한다. 나머지 방어정이 고정되어 있으므로 이 차이는 $p$ 의
선택에서만 온다. COMA~\citep{foerster2018coma}의 반사실적 기준선과 같은 착상이되 학습된
비평가 대신 시뮬레이터를 실제로 재진행하며, 그 대가로 결정당 시뮬레이터 호출이 $G$ 회에서
$GP$ 회로 는다.

\textbf{채택 여부.} 같은 예산(800 업데이트)의 스크리닝에서 반사실 모드는 600 업데이트 시점
greedy 포획률 $0.801$ 로 팀 이득 모드의 $0.800$ 과 동률이었고 시뮬레이터 호출은 3배였다.
본 논문이 평가하는 체크포인트는 \emph{팀 이득 손실}(식~\eqref{eq:lpg}--\eqref{eq:loss})로
학습한 것이다.

\subsubsection{보상 설계}
\label{sec:reward}

한 후보의 창 보상 $R(\mathbf{s}, \mathbf{a}; H)$ 는 결정 시점부터 $H$\,step 동안 일어난
일을 하나의 스칼라로 합산한 팀 공유 보상이며, 이벤트 항·퍼텐셜 차분·조밀 항으로 구성된다.
\begin{equation}
R = R_\mathrm{event} + \big[\gamma\,\Phi(\mathbf{s}_{t+H}) - \Phi(\mathbf{s}_t)\big] + R_\mathrm{dense}
\label{eq:reward}
\end{equation}
이벤트 항은 창 안의 사건 수에 고정 가중을 곱한 합이다.
\begin{equationsmall}
\begin{aligned}
R_\mathrm{event} = \;& r_{\mathrm{cap}} N_{\mathrm{cap}} + r_{\mathrm{brc}} N_{\mathrm{brc}}
  + r_{\mathrm{col}} \big(N_{\mathrm{col}} + N_{\mathrm{land}}\big)\\
  & + r_{\mathrm{tng}} N_{\mathrm{tng}}
  + \mathbb{1}_{\mathrm{end}}\big(r_{\mathrm{wipe}}\,\mathbb{1}[\mathcal{A} = \varnothing] + r_{\mathrm{surv}}\,|\mathcal{A}|\big)
\end{aligned}
\label{eq:revent}
\end{equationsmall}
계수는 포획 $r_{\mathrm{cap}} = +13.5$, 돌파 $r_{\mathrm{brc}} = -8.0$, 아군·모선·육지 충돌
$r_{\mathrm{col}} = -30$, 설치 그물 접촉 $r_{\mathrm{tng}} = -13$, 창 안에서 에피소드가 끝날
때($\mathbb{1}_{\mathrm{end}}$) 전멸 $r_{\mathrm{wipe}} = +8$, 잔존 적 1척당 $r_{\mathrm{surv}} = -5$ 이다. $\Phi$ 는 적의
모선 근접도에 대한 퍼텐셜로, 차분 형태로 주어지므로 최적 정책을
보존한다~\citep{ng1999shaping}.
\begin{equation}
\Phi(\mathbf{s}) = -w_\mathrm{threat} \sum_{m \in \mathcal{A}}
          \mathrm{clip}\!\left(1 - \frac{\lVert \mathbf{e}_m - \mathbf{m} \rVert}{R_w},\, 0,\, 1\right)
\label{eq:potential}
\end{equation}
여기서 $\mathcal{A}$ 는 생존 적선 집합, $R_w = 6300\unit{m}$ 는 교전 해역 반폭,
$w_\mathrm{threat}=3.0$, $\gamma = 0.99$ 이다. 적이 포획되어 $\mathcal{A}$ 에서 빠지거나
멀어지면 $\Phi$ 가 오른다.

조밀 항 $R_\mathrm{dense}$ 는 창 안의 기동 품질에 대한 비용·보너스의 합이다: 이동거리·소모
그물·선회량·직전 결정 대비 목표점 이동에 대한 비용, 모선·타 방어정 근접에 대한 제곱
페널티, 그리고 적$\to$모선 직선이 설치된(또는 계획된) 그물에 막힌 비율에 대한 커버리지
보너스이다. 커버리지는 창 안에 포획이 없을 때도 후보를 변별하는 선행 지표다. 이 항들의
가중은 커버리지를 제외하면 이벤트 계수의 수십분의 일 규모이고, 커버리지 항은 최대 $4.2$ 로
포획 1건($13.5$)의 3분의 1을 넘지 않게 두었다.

퍼텐셜 항은 이벤트 항의 희소성을 메운다. 216개 결정의 실측에서 포획은 평가 창의 46.8\,\%
에서만 발생했고, 나머지 창에서는 퍼텐셜 항이 유일한 큰 신호라 이것이 없으면 그룹 내 후보의
순위가 갈리지 않는다. 그 평균은 $+0.36$, 절댓값 평균은 $7.40$ 으로 창마다 크게 변동하되
전체적으로 상쇄되어, 정책은 바꾸지 않고 변별력만 제공한다는 설계 의도와 부합한다.

\subsubsection{초기화, 스케줄, 수렴 진단}
\label{sec:trainproc}

\paragraph{행동 복제 초기화}
\label{sec:bc}
그룹 상대 이득은 후보들이 갈릴 때만 신호를 주므로, 휴리스틱 행동(\S\ref{sec:heuristic})을
목표로 300 step 의 행동 복제로 정책을 부팅한다. 목표 $\omega_k(i)$ 는 원-핫이 아니라
휴리스틱 픽셀 $i^{*}_k$ 를 중심으로 유효 픽셀 위에 정규화한 가우시안 가중
($\sigma_{\mathrm{BC}} = 0.8$\,px)이며 --- 인접 픽셀은 사실상 같은 행동이라 원-핫은 기울기가
지나치게 희소하다 --- 두 번째 픽셀의 질의는 휴리스틱의 첫 픽셀로 갱신한다(teacher forcing).
\begin{equation}
\mathcal{L}_{\mathrm{BC}}(\theta) = -\frac{1}{P}\sum_{p=1}^{P} \sum_{k=0}^{K-1} \sum_{i \in \mathcal{V}^{(k)}_p}
  \omega_k(i)\, \log \pi^{(k)}_p(i)
\label{eq:lbc}
\end{equation}

\paragraph{학습률·엔트로피 계수 스케줄}
\label{sec:schedule}
옵티마이저는 Adam 이다. 학습률은 $\eta_0 = 10^{-4}$ 에서 $0.05\,\eta_0$ 까지 단일 주기
코사인 어닐링~\citep{loshchilov2017sgdr}($T = 800$ 업데이트)으로, 식~\eqref{eq:loss}의
엔트로피 계수는 $\lambda_0 = 0.01$ 에서 $0.1\,\lambda_0$ 까지 같은 예산 축 위에서 선형으로
줄인다. 두 계수를 같은 축에서 줄이면 분포가 뾰족해지는 시점과 보폭이 작아지는 시점이
맞물린다.

\begin{figure*}[tp]
\centering
\includegraphics[width=\textwidth]{fig13.png}
\caption{수렴 진단(회색 시드별, 파란 실선 시드 평균): (a) greedy 포획률, (b) 정책 엔트로피, (c) 근사 KL 발산,
         (d) 파라미터 이동}
\label{fig:convergence}
\end{figure*}

\paragraph{수렴 진단}
\label{sec:convergence}
채택 런은 20 업데이트마다의 요약만 남겼으므로, 같은 설정을 시드 3개로 다시 돌려 매
업데이트마다 수렴 진단을 기록하였다(Fig.~\ref{fig:convergence}; (a)의 빨간 파선은 채택 런,
초록 파선은 휴리스틱 기준선). 진단량은 갱신 전후 정책 사이의 근사 KL(같은 관측·같은 표본
후보에 대한 확률비의 $k_3$ 추정량)과 행동복제 초기점으로부터의 상대 파라미터 이동이다.
세 시드에서 지표들은 같은 방향을 가리킨다: 근사 KL은
처음 100 업데이트 평균 $1.4\times10^{-2}$ 에서 마지막 100 업데이트 $1.5\times10^{-4}$ 로 두
자릿수 줄고, 파라미터 이동은 뒤 절반에서 포화하며($0.031 \to 0.034$), 정책 엔트로피는
$2.42$ 에서 $2.06$\,nat 로 내려가 400 업데이트 이후 평탄하다. greedy 포획률은 세 시드 모두
800 업데이트에서 기준선 $0.749$ 위에 있고($0.817$, $0.777$, $0.801$), 채택 런의 평가점도 이
범위 안에 든다. 그래디언트 노름은 식~\eqref{eq:adv}의 이득이 표준편차로 정규화된
무차원량이라 수렴 지표로 쓰지 않는다.

\subsection{언어모델 기반 전략 계층}
\label{sec:commander}

전략 계층의 위치는 Fig.~\ref{fig:architecture}에 있다. 입력은 전장 상태·임무 목표·교전
규칙·자산 상태이고, 출력은 무리별 배정과 대기(HOLD)이며, 그 배정이 기동 계층으로 전달되어
각 방어정의 그물벽 위치 결정을 조건짓는다. Fig.~\ref{fig:runshot}은 이 계층이 실제로
동작하는 화면으로, 오른쪽 패널이 지휘관이 낸 무리별 배정과 그 근거(모델 출력을 영문으로
옮김)이고, 왼쪽이 그 배정 아래에서 방어정이 회랑에 그물벽을 친 전장이다.

\begin{figure*}[tp]
\centering
\includegraphics[width=0.9\textwidth]{fig14.png}
\caption{실행 화면(양동 포메이션, Gemini 지휘관) --- 전장(왼쪽)과 지휘관의 배정·근거(오른쪽)}
\label{fig:runshot}
\end{figure*}

언어모델에는 좌표 원시 데이터 대신 \emph{기하량이 미리 계산된} 전장 상태를 직렬화하여
전달한다. 각 적 무리의 요격점까지의 이동거리와 선회량, 모선 기준
방위, 설치된 그물에 의한 접근로 차단 여부 등이 포함된다. 언어모델이 삼각함수를 암산하게
하는 대신 판단에 필요한 양을 미리 제공하는 것이다.

\subsubsection{구조적 출력에 의한 행동 규제}
\label{sec:schema}

언어모델의 출력을 자유 문장으로 받으면 하위 계층이 파싱에 실패할 수 있고, 실패 시 그
주기의 배정이 통째로 사라진다. 본 연구는 출력을 \emph{스키마로 강제}하여 이 실패 모드를
구조적으로 차단한다~\citep{openai2026function}. 계획을 다음 자료형으로 선언하고, 그로부터 생성한
JSON 스키마를 모델 호출에 함께 넘긴다.

% 좁은 단에 맞춰 줄을 미리 끊는다. 자동 줄바꿈에 맡기면 들여쓰기가 무너져
% 자료형 구조가 안 보인다.
\begin{quote}\ttfamily\scriptsize\setlength{\parindent}{0pt}
CommanderPlan \{\\
\hspace*{1.2em} rationale: str\\
\hspace*{1.2em} deployments: List[\{\\
\hspace*{3.0em} cluster\_id: int,\\
\hspace*{3.0em} ally\_ids: List[int] \}]\\
\hspace*{1.2em} hold\_ships: List[int]\\
\}
\end{quote}

스키마는 디코딩 단계에서 문법으로 작용하여 형식에 맞지 않는 토큰이 생성되지 않게 하며,
응답은 같은 자료형으로 역직렬화하여 검증한다. 온도는 0으로 두어 같은 전장에서 같은 배정이
나오게 하는데, 이는 \S\ref{sec:paired}의 시드 대응 설계의 전제이기도 하다.

\textbf{필드 순서.} 스키마는 \emph{판단 근거 $\to$ 무리별 배정 $\to$ 대기 목록} 순으로 고정하여 배정이 근거를 조건으로 생성되게 하되, 근거 원문은 저장하지 않았으므로 그 인용 준수 여부는 검증하지 못한다(\S\ref{sec:limitations}).

\textbf{빈 배정의 의미는 프롬프트가 규정한다.} \texttt{ally\_ids} 를 비우는 것은 ``시스템이
대신 고르라''가 아니라 ``이 무리에는 아무도 보내지 않는다''를 뜻한다. 코드가 보완하지 않으므로 빈 목록은 그대로 무방비를 낳으며, 프롬프트가 이 결과를 명시하여
지휘관이 그 대가를 알고 결정하게 한다.

\textbf{권한 범위.} 언어모델이 정할 수 있는 것은 출력 스키마가 정확히 제한한다 --- 무리별
배정 항목의 필드는 \texttt{cluster\_id} 와 \texttt{ally\_ids} 둘뿐이고, 그 밖에 결정되는 것은
판단 근거와 대기 목록(\texttt{hold\_ships})뿐이다. 즉 권한은 \emph{어느 배가 어느 무리를
맡는가}와 \emph{이번 주기에 누가 멈추는가}로 닫혀 있다. 그물 전개 여부도, 벽의 방위도, 교전
반경도 출력에 없으며, 배정으로부터 요격점을 만드는 기하는 전적으로 코드다. 언어모델이 낸
배정은 코드가 보정하지 않고 그대로 쓰이므로(비워 둔 무리는 담당 없이 남고, 지목되지 않은
방어정은 대기한다), \S\ref{sec:results}의 조건 간 차이는 배정 판단 자체에 귀속된다.

그 기하는 휴리스틱 경로와 같은 규칙을 쓴다. 한 무리에 두 척 이상이 배정되면
\S\ref{sec:heuristic}의 잉여 배정과 같은 반경 층 분리($\lambda_{\mathrm{L}} = 0.18$)와 방위
분리($\pm 40^\circ$)가 작동하되, 요격점이 요격 구역을 벗어나 유효 마스크가 붕괴하지 않도록
반경을 구역 안으로 자른다.

배정이 확정되면 각 방어정의 요격점이 정해지고, 그 요격점을 중심으로 기동 계층이 그물벽을 놓을
수 있는 후보 영역(유효 마스크)이 함께 결정된다. 즉 배정은 기동 계층의 선행 조건이며,
배정이 바뀌면 후보 영역 자체가 바뀐다.

